\section{Cyclic decomposition}

Let $V_A$ denote $\Q^n$ with the $\Q[X]$ action obtained by evaluating
polynomials at $A$.  In Lean this is mathlib's \code{Module.AEval'} at the
manifest-pinned source commit \citep{mathlib4320aeval}.
The cokernel of $XI-A$ is first related to $V_A$ by explicit mutually inverse
linear maps.  Composing this equivalence with the Smith cokernel
decomposition yields
\[
  V_A \simeq_{\Q[X]}
  \prod_{i:\operatorname{Fin} n}\Q[X]/(d_i).
\]
The product retains $\Q[X]/(1)$ terms.  These are zero modules, not missing
data.  The separate display reader removes them without changing the
equivalence.

Each nonzero monic $d_i$ gives a power basis
\[
  1,\bar X,\ldots,\bar X^{\deg(d_i)-1}
\]
of its quotient.  Pulling the product of these bases back through the cyclic
decomposition constructs a basis of $\Q^n$.  This step needs the dimension
identity
\[
  \sum_i \deg(d_i)=n.
\]
Rather than using this equality only for arithmetic, the implementation turns
it into an equivalence
\[
  \left(\sum_i \operatorname{Fin}(\deg d_i)\right)
  \simeq \operatorname{Fin} n.
\]
It therefore fixes the order and type of every coordinate in the final block
matrix.

The annihilator of the product module is the intersection of the ideals
$(d_i)$.  Because the Smith diagonal is a divisibility chain, this
intersection is generated by the last factor in positive dimension.
Mathlib identifies the annihilator of \code{Module.AEval'} with the ideal
generated by the minimal polynomial.  Both generators are monic, so
association again strengthens to equality:
\[
  d_{n-1}=\mu_A.
\]
The public option theorem states this result for every positive dimension
while preserving \code{none} at dimension zero.
